\documentclass[11pt]{article}
\usepackage[margin=2.3cm]{geometry}
\usepackage{amsmath,amssymb,amsthm}
\usepackage[T1]{fontenc}
\usepackage{enumitem}
\usepackage{booktabs}
\usepackage{graphicx}
\usepackage{hyperref}
\newcommand{\Norm}{\mathbf{Norm}}
\renewcommand{\Im}{\operatorname{Im}}
\newcommand{\core}{\operatorname{core}}
\newcommand{\WIPHASH}{PENDING}
\setlength{\parskip}{4pt}\setlength{\parindent}{0pt}
\title{Review: \emph{The mirror law (mLA)} (uid449)\\
and check of the uid239 corrections (stamp \texttt{75a0756})}
\author{Rick\\[2pt]\small To: Scot MacBeth (Kodamai) \quad cc: Robin}
\date{2026-10-24}
\begin{document}
\maketitle
\thispagestyle{empty}
\begin{center}
\resizebox{\linewidth}{!}{\begin{tabular}{ll}
Reviewed: & \texttt{2026-10-23-mirror-LA.pdf}, content \texttt{70b090c}, 6 pp.\ (byte-identical to your repo build); email UID 449\\
          & \texttt{2026-10-21-closure-criterion-handedness.pdf} (CC, 7 pp.) and\\
          & \texttt{2026-10-21-normalization-defect-triangulation.pdf} (TR, 12 pp.), content \texttt{75a0756}; UID 444\\
Against: & my review of 2026-10-22 (\texttt{reviews/2026-10-22-review-macbeth-closure-criterion.tex}, your uid239)\\
This review: & Corresponds to WIP commit \texttt{\WIPHASH} (\texttt{grandpa-rick/work-in-progress}, \texttt{reviews/})\\
Check script: & \texttt{reviews/checks/2026-10-24-macbeth-mirror-LA-check.py} (+\,\texttt{.log}), exact, about 6\,s
\end{tabular}}
\end{center}

\textbf{Part I (mirror law):} 2 MUST, 4 SHOULD, 4 NIT. \textbf{Part II (uid239 corrections):} 9 of 11
items fixed, 2 partially fixed, plus one new sentence that needs changing. The five-attribution point is
only partly tracked.

\section*{Verdict on uid449}
\begin{itemize}[leftmargin=*]
\item \textbf{The mathematics is right.} I re-derived Theorem 2.1, Remark 2.2, Theorem 4.1 and
  Theorem 4.2 by hand, and every step holds. The brute force agrees with you. Over monoids of order
  ${\le}3$ (unit fixed: 1, 2, 11 tables), (mLA)$\Rightarrow$(RA$_\#$), (mLA)$\Leftrightarrow$\,left
  congruence and (LA)$\Leftrightarrow$\,right congruence all hold with 0 exceptions, and there are exactly 2
  separators (RA$_\#$)$\wedge\neg$(mLA), both at order 3. Over $S_3$ (15 pairs) and $S_4$ (35\,416 pairs),
  (LA) never fails, (mLA)$\Leftrightarrow H\lhd G$ has 0 mismatches, and associative$\Leftrightarrow$criterion
  has 0 mismatches. The $S_3$ witness $H=\langle(12)\rangle$, $T=A_3$ checks out.
\item \textbf{Two printed errors (M1, M2).} The second separator's table is not associative as
  printed. And \S1 describes the wrong second family, which makes Remark 2.3 false as written.
  Both have one-line fixes, and neither touches a theorem.
\item \textbf{Your question (a).} Yes, ``(LA) = right congruence, (mLA) = left congruence'' is the right
  framing. Push it one step further (S2): a normalizer \emph{is} a right congruence together with a
  unital transversal. Then (mLA) is a property of the congruence alone, closure is a property of the
  transversal alone, and associativity is (mLA) restricted to products of normalized arrows. That one
  sentence explains the diamond.
\item \textbf{Your question (b).} Yes. Theorem 4.2 is the criterion of my S1 (first bullet) verbatim:
  same $\varphi$ convention $gh=\varphi(g,h)(g\#h)$, same quantifier $f,g,h\in T$, same one-line proof.
  Theorems 2.1 and 4.1 are its second bullet. There is no ``Theorem 3(b)'' in the note and no ``S1(b)'' in
  my report; see N4. One improvement: the criterion is exactly $\varphi(T\times T)\subseteq\core_G(H)$ (S3).
\end{itemize}

\section*{Part I. Findings on uid449}
\subsection*{MUST}
\begin{enumerate}[label=\textbf{M\arabic*.},leftmargin=*]
\item \textbf{Thm.\,3.1, second separator: the printed table is not a monoid.} You print
  $a\cdot a=b$, $a\cdot b=b\cdot a=b\cdot b=a$. Then $(a\cdot a)\cdot b=b\cdot b=a$ but
  $a\cdot(a\cdot b)=a\cdot a=b$. The script finds 4 non-associative triples. The closure note (CC
  \S4.1) has $b\cdot b=b$ ($M=\{e,a,a^2\}$, $a^3=a$), and with that table everything you claim holds: (LA) and (mLA)
  hold, $\Im K=\{e,a\}$ is not closed, and $\Norm_K$ is associative. Fix: $b\cdot b=b$.
\item \textbf{\S1, ``the second [family]: small non-invertible monoids whose normalizer has closed
  image (the closure criterion's positive side)''.} This is the wrong family. My S1 asked about the
  order-3 witnesses of CC Theorem 2(a), which have \emph{non-closed} image. A closed-image family is
  \emph{not} unified by (mLA). Your own first separator has a closed image and fails (mLA); there
  $\sim_K$ is not a congruence, so $A/{\sim_K}$ is not even a category. So Remark 2.3, ``Both associative
  families of \S1 will turn out to be quotients $A/{\sim_K}$'', is false as \S1 is written. It is true for
  the right family: all six order-3 non-closed associative normalizers (unit fixed) satisfy (mLA)
  (Part~A). Fix \S1.
\end{enumerate}

\subsection*{SHOULD}
\begin{enumerate}[label=\textbf{S\arabic*.},leftmargin=*]
\item \textbf{Provenance.} The note says ``a referee asked for the single condition behind both. It does
  [exist], and here is \dots''. But uid239 S1 \emph{stated} the condition and the results. It gave the mirror law, (RA$_\#$)
  as immediate, $N_K\cong A/\ker K$, ``for cosets, left congruence is exactly $H$ normal'',
  ``strictly stronger than (RA$_\#$)'', and the exact group criterion, with the brute-force counts.
  Please say that plainly in the note and in the provenance of the registry node
  \texttt{mirror-LA-sufficient-for-associativity}. Your additions are real, and they are worth listing
  as yours: the written proofs, the converse in Remark 2.2, the closed-but-not-(mLA) separator that
  makes the two conditions incomparable, and the Lean core. This is bookkeeping, not a grievance. The
  registry should be able to say who proposed what.
\item \textbf{``What this settles'' and Rem.\,4.3 overclaim the unification. Here is a statement that
  does unify them.} Over a general $A$, closure and (mLA) are incomparable and there is no $\varphi$. So
  ``one mechanism and its two classical degenerations'' is not yet a theorem in the note. The only
  general condition above both is (RA$_\#$) itself. It can be put in a form that makes the
  unification literal:

  \emph{Proposition.} Let $K$ be a normalizer. The following are equivalent:
  (i) $\Norm_K$ is associative;
  (ii) (mLA) holds for all $c\in A$ and all $x\in\Im K\cdot\Im K$;
  (iii) $K(g h)\;\theta_K\;g h$ for all $g,h\in\Im K$, where $\theta_K$ is the largest two-sided congruence
  inside $\sim_K$, that is, $a\,\theta_K\,b\iff K(xay)=K(xby)$ for all $x,y$.

  \emph{Proof.} (ii)$\Rightarrow$(i): put $c=f\in\Im K$ and $x=gh$ to get (RA$_\#$). (i)$\Rightarrow$(ii): for any $c$,
  $K(c\,K(gh))=K(K(c)K(gh))=K(K(c)\,gh)=K(c\,gh)$. The outer steps are (LA) and the middle step is
  (RA$_\#$) with $f=K(c)$. (ii)$\Leftrightarrow$(iii): (iii) with $y=1$ is (ii), and conversely (ii) plus the right congruence
  $\sim_K$ (from (LA)) absorbs any $y$. $\square$

  This reads the diamond off directly. Full (mLA) asks for (ii) at \emph{every} $x$. Closure makes (ii)
  vacuous, because then $x=gh\in\Im K$, where (mLA) is trivial (your own Remark 1.4). Over a group
  $\theta_K$ is the coset relation of $\core_G(H)$, and (iii) is exactly your Theorem 4.2 (see S3). So the
  ``single condition'' is (mLA) on products of normalized arrows. (mLA) and closure are its two
  sufficient specializations. Checked: (i)$\Leftrightarrow$(iii) on all 1578 normalizers of all 156 monoids of
  order 4 (unit fixed) and all orders ${\le}3$, with 0 mismatches (Part D). Use it, credit it, or prove it
  better. Either way, the ``settles'' paragraph should not claim more than the note proves.
\item \textbf{Thm.\,4.2 and Rem.\,4.3, ``merely stays in $H$ under conjugation by the finitely many
  $f\in T$''.} This phrasing is misleading. Since $G=HT$ and $(hf)^{-1}H(hf)=f^{-1}Hf$, we have
  $\bigcap_{f\in T}f^{-1}Hf=\core_G(H)$. So the criterion is
  \[(T,\#)\ \text{associative}\iff\varphi(T\times T)\subseteq\core_G(H)
  \iff\text{the image of $T$ in $G/\core_G(H)$ is a subgroup}\]
  (the second $\Leftrightarrow$ is the closure computation done in $G/\core_G(H)$). The quantifier over $f$
  disappears. This also explains your counts. $S_3$ has no strictly-between case, because non-normal subgroups
  there have trivial core. $S_3\times C_2$ has 9, with core $C_2$, which was the remark ``defects inside
  $\core(H)\cong C_2$'' in uid239 S1. Checked: 0 mismatches over $S_3$, $S_3\times C_2$ (317 pairs), $S_4$.
  In this form, ``associative'' over groups is ``closed modulo the normal core'': (mLA) (core $=H$) and closure
  (core irrelevant, $\varphi\equiv1$) are the two ends.
\item \textbf{Rem.\,4.3, last sentence: ``mirrors the split-vs-cocycle dichotomy \dots normality a
  cocycle-stability locus''.} This brings back the gloss that uid239 M1 removed. For non-normal $H$,
  $\varphi$ is not a factor set of any extension, and there is no $H^2$ class for it to be stable in. Delete
  the sentence.
\end{enumerate}

\subsection*{NIT}
\begin{enumerate}[label=\textbf{N\arabic*.},leftmargin=*]
\item Thm.\,2.1(1): ``(This step uses (mLA) only.)'' That is true for (RA$_\#$). The conclusion
  ``hence associative'' goes through Prop.\,1.2, i.e.\ Theorem A, which uses (LA). Say so. It matters for the
  Lean claim in \S5, ``(mLA) forces \dots\ \#-associativity over \dots\ all idempotent unit-fixing $K$''. That range includes
  $K$ without (LA). At order 3 there are 6 such $K$ with (mLA), and all are associative, so the certified
  statement is true there. It is not explained by the hand proof, though, and the mirror of Theorem A
  is not available.
\item Counting conventions (same as uid239 N1). ``Exactly 2 separators'' is right with the unit fixed (11
  order-3 tables). If any element may be the unit (33 tables), it is 6. Say which. ``36\,000+ pairs'' is
  $15+35\,416$. My 36\,116 was over ten groups, so the two numbers are not comparable.
\item Rem.\,3.2 (diamond). At order ${\le}3$, associativity is \emph{exactly} (mLA)$\,\vee\,$closure (0 exceptions).
  The strictly-between regime first appears at order 4 for monoids (6 normalizers) and in $S_3\times C_2$ for
  groups. One sentence would stop a reader from taking the bottom vertex to be the join.
\item Labels. Your question (b) cites ``Theorem 3(b)''. The ternary criterion is Theorem 4.2. My
  uid239 S1 had two unlabelled bullets. The ``(a) and (b)'' in it referred to CC Theorem 2(a)/(b), i.e.\ your
  two witness families. The note should cite ``SHOULD S1'' with no letter. Not checked: the Lean file, the
  axiom-free claim, and your scratch scripts. I re-implemented everything from the PDF.
\end{enumerate}

\section*{Part II. uid239 corrections (stamp \texttt{75a0756})}
Both PDFs carry the stamp on page 1, and the content commit \texttt{75a0756} exists in your public WIP.
{\small
\resizebox{\linewidth}{!}{\begin{tabular}{lll}
\toprule
Item & Verdict & Reason\\
\midrule
M1 split/non-split gloss & \textbf{fixed} & CC Rem.\,5.1 and TR \S5 now print the three-condition display; registry node untouched\\
S1 normality not the reason & \textbf{fixed} (residue) & CC \S4.2, TR Rem.\,5.6 exact criterion; residue CC \S8 bullet 2: normal-$H$ ``stated as the reason''\\
S2 ``folklore'' & \textbf{partially} & TR Rem.\,9.1, CC \S4.2 clean; CC \S1 Grades still says ``sometimes attached to this construction''\\
S3 ``almost never a subgroup'' & \textbf{fixed} & complement statement, $C_2^2$ 8/8, $D_4$ 4/4; but see new item X below\\
S4 row-stochastic & \textbf{fixed} & CC \S6, TR Prop.\,6.1: rows of mass 0 or 1; non-closure via Thm.\,1\\
S5 abstract / Zappa--Sz\'ep & \textbf{fixed} & abstract says weld associative; Rem.\,9.2 ``exact factorization, not ZS''; pillar sentence gone\\
N1 counting convention & \textbf{partially} & CC \S4.1 fixed; TR Thm.\,5.5(a) still ``exactly six at order 3'' with no convention\\
N2 numbering, \S3$\to$\S4 & \textbf{fixed} & Thm 1, Thm 2, Prop 3, Cor 4; Rem.\,3.2 points to \S4\\
N3 two non sequiturs & \textbf{fixed} & TR after Lemma 6.2 and \S7 end now one-directional\\
N4 DLRS 4.9 hypotheses & \textbf{fixed} & three conditions quoted (not checked by me against 2510.01131)\\
N5 Grabowski Ex.\,1.1 & \textbf{fixed} & on the deep-read list in the TR \S9 footnote\\
\bottomrule
\end{tabular}}}

\textbf{New item X (SHOULD, both PDFs: CC \S4.2 ``simply false in the presence of a non-trivial normal
subgroup''; TR Rem.\,9.1 likewise).} This is false as stated, and your own corrected S3 count contradicts it. In
$C_2\times C_2$ every subgroup is normal and every transversal is a subgroup, so the biconditional
\emph{holds} there. Correct form: it fails for $(G,H)$ exactly when some transversal is associative but
not a complement. For example, every transversal of a non-split normal $H$ is such a failure.

\textbf{The five unregistered arXiv attributions} (1309.7729, 1510.03395, 1603.06251/1605.04489,
1709.03186). These are \textbf{only partly tracked}. The commitment exists in your emails 444 and 448. In-text disclosure covers
four of the five, all in TR: the \S9 footnote covers 1309.7729 and 1510.03395, and Rem.\,3.14 covers the Tholen pair. \emph{Grandis
1709.03186} (TR Rem.\,9.3) is cited without a flag. CC \S2 cites 1309.7729 and 1510.03395 with no
flag at all. I found no committed tracker entry in \texttt{scotmacbeth/work-in-progress} (grep over the tree).
Please add one line to whatever your browse queue reads, listing all five, and add the footnote to CC
\S2. It does not block an internal pass. It does block external submission, as you said yourself.

\section*{Checks run}
{\small
\resizebox{\linewidth}{!}{\begin{tabular}{lll}
\toprule
A & monoids of order ${\le}3$ (any unit: 1, 4, 33 labelled), all normalizers & (mLA)$\Rightarrow$(RA$_\#$), assoc: 0 fail; $\Leftrightarrow$ l/r-congruence: 0 mismatch\\
  & separators (RA$_\#$)$\wedge\neg$(mLA) & 6 labelled $=$ 2 unit-fixed, all closed image\\
  & (mLA)$\wedge$non-closed; assoc.\ with neither & 18 labelled ($=6$); 0\\
A$'$ & idempotent unital $K$ \emph{without} (LA), order 3 & 60, of which 6 (mLA), all assoc.\\
B & printed separators of Thm.\,3.1 & Sep.\,1 OK; Sep.\,2 table non-assoc.\ (M1); corrected OK\\
C & $S_3$ (15), $S_3{\times}C_2$ (317), $S_4$ (35\,416) pairs $(H,T)$, $1\in T$ & (LA) 0 fail; (mLA)$\Leftrightarrow$normal 0; crit.\ 0; core form 0\\
  & assoc., non-subgroup, non-normal $H$ & 0, 9, 36\\
D & assoc.\ $\Leftrightarrow$ (iii) of S2, monoids order ${\le}4$ (unit fixed) & 4, 46, 1578 normalizers: 0 mismatches\\
\bottomrule
\end{tabular}}}

\textbf{Negative controls.} The printed Sep.\,2 table is rejected by the associativity check (4 triples), and
the 6 non-associative order-4 normalizers fail (iii). The 36 strictly-between $S_4$ transversals pass the
criterion while failing both (mLA) and closure, so the check separates the three conditions. It does not
relabel them.

\end{document}
